\begin{lemma}[An extending gluing bundle]
\hypertarget{tag-03I9}{}\label{wbII:lem:extending-gluing}
\StatementRecord{wbII:lem:extending-gluing}{proved}{03I9}{\SeriesRef{cdp:st:B1}, \SeriesRef{wbII:lem:gluing-bundle}}

After shrinking $\Delta$, the sections
\[
 O(t)=[0:1:0],\qquad
 P(t)=[3:\sqrt{36+t}:1]
\tag{W.162}\hypertarget{tag-03E0}{}\label{wbII:eq:CN:4.5a}
\]
of $\mathcal C\to\Delta$ are disjoint and holomorphic.  If
$B=\cO_{\mathcal C}(P-O)$, then the line bundle
$\widetilde A=\operatorname{pr}_1^*B$ on $\mathcal X=\mathcal C\times\PP^1$
restricts to $\operatorname{pr}_1^*A_{\lambda}$ on $S$ for
$\lambda=1/3$ (or its inverse, according to the gluing convention).  The bundle in
\cref{cdp:st:B1} with this value of $\lambda$ therefore extends to
$\mathcal X$.
\end{lemma}



\begin{proof}
The equation of $\mathcal C$ is satisfied by both sections; the square root
is holomorphic for $|t|$ small, and the two sections are disjoint because
$O$ has $z=0$ whereas $P$ has $z=1$.  On the normalization of $C_0$, write
$r=u/v$.  Then $O(0)$ corresponds to $r=\infty$, $P(0)$ to $r=2$, and the
two branches over the node correspond to $r=1$ and $r=-1$.  Hence
\[
 \nu^*(B|_{C_0})\simeq\cO_{\PP^1}([2]-[\infty]).
\]
The rational function $f(r)=r-2$ has precisely this divisor and trivializes
the pullback.  Its values at the two branches are $f(1)=-1$ and
$f(-1)=-3$, so the induced gluing multiplier is $f(1)/f(-1)=1/3$,
up to inversion.  This is not a root of unity, so its gluing bundle is non-torsion by
Lemma~\ref{wbII:lem:gluing-bundle}.

\end{proof}